\chapter{The Multi-Manager Platform}\label{ch:fm:the-multi-manager-platform}

A large multi-manager platform's regulatory brochure lists, among the expenses its funds bear, the salaries, bonuses and performance-based
compensation of the portfolio managers who run its teams, their recruiters' fees, their computers and their offices. It calls the arrangement an
``expense pass-through''. It adds that a portfolio manager is paid a share of the profit of his own team, ``without taking into account the
performance of other Portfolio Managers or of the Millennium Partners Funds generally'', and that portfolio managers with positive performance
``will receive performance-based compensation even if the Millennium Partners Funds' overall returns are negative''. The investor pays each
winning team's bonus, including in the years the fund loses. This chapter prices that promise.

\section{Pods, the centre and the platform's promise}

A multi-manager platform (One Quant Book 1, chapter 1) is a hedge fund run as many independent teams, each a pod (One Quant Book 8, chapter 28),
under a centre that allocates capital, sets risk limits, nets and hedges the sum, and pays for the infrastructure. Its promise to investors is
diversification of talent: fifty teams with modest Sharpe ratios and low correlation add up to a fund with a high one. Its promise to portfolio
managers is a formulaic share of what their own team makes, paid whatever the others do.

\begin{proposition}[Diversification of talent]\label{prop:fm:the-multi-manager-platform:div}
Let $n$ pods have the same volatility $\sigma$ and Sharpe ratio $S$ and pairwise correlation $\rho$, and let the fund hold them in equal weights.
The fund's Sharpe ratio is
\[
S_{\mathrm{fund}}=S\sqrt{\frac{n}{1+(n-1)\rho}}\;\longrightarrow\;\frac{S}{\sqrt\rho}\quad(n\to\infty).
\]
\end{proposition}

\begin{proof}
The fund's mean is $S\sigma$ and its variance $\sigma^2(1+(n-1)\rho)/n$ per unit of capital split equally; divide. The limit follows as the
$(n-1)\rho$ term dominates.
\end{proof}

\begin{center}
\small
\begin{tabular}{@{}lrrrr@{}}
\toprule
fund Sharpe ratio (pods at $S=1$) & $n=10$ & $n=25$ & $n=50$ & $n=100$\\
\midrule
$\rho=0.05$ & 2.63 & 3.37 & 3.81 & 4.10\\
$\rho=0.10$ & 2.29 & 2.71 & 2.91 & 3.03\\
$\rho=0.20$ & 1.89 & 2.08 & 2.15 & 2.19\\
\bottomrule
\end{tabular}
\end{center}

The table is the platform's business case and its limit. Correlation, not the number of teams, sets the ceiling: at $\rho=0.1$ it is
$1/\sqrt{0.1}=3.16$ times a pod's Sharpe ratio, and fifty pods already reach 2.91 of it. Past a few dozen teams, adding pods buys little unless
they are less correlated with the ones already there; lowering correlation, by hiring different strategies or hedging what they share
(\cref{def:fm:the-multi-manager-platform:centre}), buys more than hiring.

\begin{omfigure}
\begin{tikzpicture}[font=\footnotesize,>=stealth,box/.style={draw,rounded corners,minimum height=0.62cm,align=center,fill=omDef!10}]
\node[box,minimum width=2.4cm] (inv) at (0,3.6) {investors};
\node[box,minimum width=3.0cm] (fund) at (0,2.2) {the fund\\{\scriptsize capital, costs, fees}};
\node[box,fill=omProp!12,minimum width=3.0cm] (centre) at (5.2,2.2) {the centre\\{\scriptsize allocation, risk, centre book}};
\foreach \x/\n in {-3.6/1,-1.8/2,0/3,1.8/4,3.6/5} {\node[box,fill=gray!12,minimum width=1.5cm] (p\n) at (\x,0.4) {pod \n};}
\node[font=\scriptsize] at (5.2,0.4) {\dots\ fifty pods};
\draw[->] (inv) -- node[right,font=\scriptsize] {capital} (fund);
\foreach \n in {1,2,3,4,5} {\draw[->] (fund.south) -- (p\n.north);}
\draw[<->,dashed] (fund.east) -- (centre.west);
\draw[->,omThm] (p1.south) -- ++(0,-0.5) -| node[pos=0.25,below,font=\scriptsize] {each team paid a share of its own profit} (p5.south);
\node[font=\scriptsize,omThm,align=left,anchor=west] at (-6.0,2.2) {pass-through:\\all teams' costs\\and pay};
\end{tikzpicture}

\omcaption{A multi-manager platform. Investors' capital is allocated to pods by the centre, which also hedges or adds to their netted exposures
through its own book; the fund bears every team's costs and pay (pass-through) and each team is paid on its own profit. Model:
\texttt{firm.podshop}.}\label{fig:fm:the-multi-manager-platform:structure}
\end{omfigure}

\begin{definition}[Pass-through fee]\label{def:fm:the-multi-manager-platform:pass}
A \emph{pass-through fee}\index{pass-through fee} is a fund's charge to its investors of the actual costs of running it, the pay of its
investment teams and staff, its data, systems and offices, instead of, or in addition to, a fixed management fee out of which the manager would
pay them.
\end{definition}

\begin{definition}[Payout rate]\label{def:fm:the-multi-manager-platform:payout}
A team's \emph{payout rate}\index{payout rate} is the share of its own trading profit, net of its costs and of any losses it carries forward, that
is paid to the team as performance compensation.
\end{definition}

\begin{dated}{2026-09}{What one platform's brochure says}\label{dat:fm:the-multi-manager-platform:adv}
Millennium Management LLC, Form ADV Part 2A (25 August 2026): its funds ``generally bear, directly or indirectly, all expenses incurred in
connection with'' their operation, an arrangement it calls an ``expense pass-through'', including the compensation of portfolio managers and other
employees, recruiters' fees, equipment and offices; this is ``separate from and in addition to'' the performance- and asset-based compensation the
adviser receives, and is deducted before it. A portfolio manager's compensation is ``generally determined as a percentage of profits earned by
such Portfolio Manager during the preceding calendar year''; losses are carried forward; there is ``generally no carryback or clawback''. The
adviser allocates and reallocates capital among portfolio managers and makes direct investments of the funds' capital, including ``hedges, or
contra trades that seek to establish a reduction in certain exposures'' and trades that increase exposure ``to netted positions held by a number
of Portfolio Managers''.
\end{dated}

\section{Pass-through fees and the investor's bill}

The chapter's platform (illustrative) runs 50 pods, each with a Sharpe ratio of 1.0 and a volatility of 20\% a year on its share of the fund's
capital, correlated at 0.1 through one shared factor. Teams are paid 20\% of their own profit, losses carried forward; a team that loses more
than 10\% of its capital in a year is replaced by a new one. Pods' costs of 3\% a year and the platform's 1\% pass through; the manager takes a
20\% performance allocation on what is left. Over twenty simulated years and twenty seeds, a hundred dollars of capital earns this a year
(\cref{fig:fm:the-multi-manager-platform:waterfall}): \$20.09 gross, of which teams receive \$4.27, costs take \$4.00 and the manager \$2.38,
leaving the investor \$9.44. The fund lost money in 9 of its 400 years.

\begin{omfigure}
\begin{tikzpicture}
\begin{axis}[width=12cm,height=5.2cm,ybar stacked,bar width=22pt,ymin=0,ymax=22,ylabel={\small \$ per \$100 of capital a year},
  xtick={0,1,2,3,4},xticklabels={gross,{team payouts},{pass-through costs},manager,investor},
  xticklabel style={font=\scriptsize,text width=1.5cm,align=center},tick label style={font=\footnotesize},
  table/col sep=comma,enlarge x limits=0.12,grid=major,grid style={gray!20},axis x line*=bottom]
\addplot[draw=none,fill=none] table[x=k,y=bottom]{figdata/desk/03-the-multi-manager-platform/waterfall.csv};
\addplot[fill=omThm!50,draw=omThm] table[x=k,y=cost]{figdata/desk/03-the-multi-manager-platform/waterfall.csv};
\addplot[fill=omDef!55,draw=omDef] table[x=k,y=total]{figdata/desk/03-the-multi-manager-platform/waterfall.csv};
\end{axis}
\end{tikzpicture}

\omcaption{Where a platform's gross profit goes, per \$100 of capital a year: team payouts on each pod's own profit, pass-through costs, the
manager's performance allocation, and what the investor keeps. Fifty pods of Sharpe ratio 1.0, twenty years, twenty seeds; illustrative terms.
Data: \texttt{fm\_platform.waterfall}.}\label{fig:fm:the-multi-manager-platform:waterfall}
\end{omfigure}

A classic fund charging 2\% of capital and 20\% of profit on the same pods would leave its investors \$14.47, charging them \$5.62. It could not
exist: out of \$5.62 of fees its manager would have to pay the teams \$4.27 and the costs \$4.00, losing \$2.65 on every hundred dollars. The
pass-through is the fee structure a platform needs once talent and infrastructure cost more than a fixed fee can pay; the investor's question is
whether what is left after them, \$9.44 here, is worth the fund's risk.

\omcode{code/firm/podshop/firm_podshop.py}{48}{66}{Each team's payout on its own profit, with losses carried forward and teams replaced; the
netting cost is what the payouts cost beyond the same rate on the fund's total.}\label{lst:fm:the-multi-manager-platform:payouts}

\section{Paying winners: the netting problem}

\begin{definition}[Netting risk]\label{def:fm:the-multi-manager-platform:netting}
A multi-manager fund's \emph{netting risk}\index{netting risk} is the cost to its investors of paying each team a share of its own profit
without netting it against other teams' losses: the difference between the sum of the teams' payouts and the same \omterm{def:fm:the-multi-manager-platform:payout}{payout rate} applied to the
fund's total profit.
\end{definition}

\begin{proposition}[The price of paying winners]\label{prop:fm:the-multi-manager-platform:netting}
Let a pod's annual profit be $X\sim\mathcal N(S\sigma,\sigma^2)$, with Sharpe ratio $S$, and pay a rate $p$ on $\max(X,0)$ each year with no loss
carried forward. If the fund's total profit is positive, the netting cost per pod is
\[
p\,\bigl(\E[\max(X,0)]-\E[X]\bigr)=p\,\sigma\,\bigl(\varphi(S)-S\,\Phi(-S)\bigr),
\]
a share $p\,(\varphi(S)-S\Phi(-S))/S$ of the pod's expected profit.
\end{proposition}

\begin{proof}
$\E[\max(X,0)]=\mu\Phi(\mu/\sigma)+\sigma\varphi(\mu/\sigma)$ for $X\sim\mathcal N(\mu,\sigma^2)$; subtract $\mu=S\sigma$ and use
$1-\Phi(S)=\Phi(-S)$. When the fund's total is positive the payout on the total is $p$ times the sum of the $X$'s, whose expectation is $p$ times
the sum of the $\E[X]$.
\end{proof}

At $p=20\%$ the netting cost is 1.7\% of gross profit for pods of Sharpe ratio 1, 7.9\% at 0.5 and 22.9\% at 0.25
(\cref{fig:fm:the-multi-manager-platform:netting}). Carrying losses forward removes most of it for a team that stays: in the chapter's platform
with no team ever replaced, netting costs 0.1\% of gross profit. Replacing losing teams restores it, because a replaced team's carried loss dies
with it and its successor starts at zero: with teams replaced after a 10\% loss, the simulated cost is 1.3\% of gross profit at a Sharpe ratio of
1 and 6.2\% at 0.5, close to the formula's bound. A platform's \omterm{def:fm:the-multi-manager-platform:netting}{netting risk} is therefore set by the quality of its teams and the rate at which it
turns them over, not by the \omterm{def:fm:the-multi-manager-platform:payout}{payout rate} alone.

\begin{omfigure}
\begin{tikzpicture}
\begin{axis}[width=10.5cm,height=5.2cm,xlabel={\small Sharpe ratio of each pod},ylabel={\small netting cost (\% of gross profit)},
  xmin=0.2,xmax=2.05,ymin=0,ymax=25,tick label style={font=\footnotesize},grid=major,grid style={gray!20},table/col sep=comma,
  legend style={font=\scriptsize,at={(0.5,-0.3)},anchor=north,/tikz/every even column/.append style={column sep=8pt}},legend columns=2]
\addplot[omThm,thick,mark=*] table[x=sr,y=closed]{figdata/desk/03-the-multi-manager-platform/netting.csv};
\addplot[omDef,thick,dashed,mark=square*] table[x=sr,y=simulated]{figdata/desk/03-the-multi-manager-platform/netting.csv};
\legend{no carry-forward (\cref{prop:fm:the-multi-manager-platform:netting}),carry-forward and teams replaced after a 10\% loss}
\end{axis}
\end{tikzpicture}

\omcaption{The cost of paying each team on its own profit, beyond the same 20\% rate on the fund's total, as a share of gross profit, against the
pods' Sharpe ratio: the closed form without carry-forward, and the chapter's platform, fifty pods correlated at 0.1, twenty years, twenty seeds.
Data: \texttt{fm\_platform.netting\_vs\_sr}.}\label{fig:fm:the-multi-manager-platform:netting}
\end{omfigure}

\begin{example}[Two pods]\label{ex:fm:the-multi-manager-platform:two}
Two pods of \$100 million each; one makes \$10 million, the other loses \$10 million. The fund's profit is zero. At a 20\% \omterm{def:fm:the-multi-manager-platform:payout}{payout rate} the first
team receives \$2 million and the second nothing: the investors pay \$2 million, plus both pods' costs, for a year in which the fund made nothing.
If the losing team stays, its first \$10 million of later profit earns it nothing; if it is replaced, its successor is paid from its first dollar.
\end{example}

\section{Capital allocation and drawdown rules}

The centre allocates capital to teams and takes it back. One Quant Book 8, chapter 28 simulated the central trade-off: a drawdown limit
(defined there) stops good teams in bad luck and lets dead ones run, and its rate of false stops is set by the limit in units of the team's own
volatility. Platforms describe their rules in general terms; the specific thresholds of named firms are not in their public documents, and this
book does not repeat second-hand figures. The model therefore uses an illustrative two-step ladder: a team's capital is halved after a 10\%
drawdown within a year and withdrawn after 15\%, and restored at the next year's start.

\begin{method}[Evaluating a platform's drawdown ladder]\label{met:fm:the-multi-manager-platform:ladder}
\begin{enumerate}
\item Express each threshold in units of the team's volatility on its capital (a 10\% cut is half a year's volatility at 20\%).
\item Simulate the platform's pods with and without the ladder, on the same random draws.
\item Report the share of team-years cut and stopped, and the fund's return, volatility and Sharpe ratio in both cases.
\item Report separately what the ladder does to dead teams (chapter 8): a ladder that improves nothing in the good-team world must earn its cost
in the bad-team one.
\end{enumerate}
\end{method}

In the chapter's platform the ladder cuts 88.4\% of team-years and stops 26.6\%: with a volatility of 20\%, a 10\% drawdown within a year is
ordinary. The fund's return falls from 20.1\% to 13.4\% and its volatility from 6.9\% to 4.8\%; its Sharpe ratio falls from 2.92 to 2.77. Every
pod in this model has a live edge, so the ladder can only cost; chapter 8 measures what it buys when some edges die.

\section{Centre books and overlays}

\begin{definition}[Centre book, factor overlay]\label{def:fm:the-multi-manager-platform:centre}
A platform's \emph{centre book}\index{centre book} is the book the centre trades itself with the fund's capital, alongside the pods: to hedge
exposures the pods share, to add to positions many pods hold, or to trade its own strategies. A \emph{factor overlay}\index{factor overlay} is a
centre-book position that offsets the sum of the pods' exposures to a common factor, so that the fund keeps the pods' specific returns without
their shared risk.
\end{definition}

Fifty pods correlated at 0.1 give a fund whose volatility is dominated by what they share: $\sigma\sqrt{(1+49\rho)/50}$ is 6.9\% for pods of 20\%
volatility, against 2.7\% if the shared part is hedged. In the model the overlay leaves the fund's mean return at 20.1\% and lifts its Sharpe
ratio from 2.9 to 7.5 (\cref{fig:fm:the-multi-manager-platform:overlay}). That is an upper bound: here all the shared risk is one factor the
centre can trade. Real pods share crowded positions and liquidity as well as market and style factors (One Quant Book 7, chapter 28), and only
the tradable part can be hedged; the brochure of the hook describes both uses of the centre's own trades, hedging netted exposures and adding to
them.

\begin{omfigure}
\begin{tikzpicture}
\begin{axis}[width=10.5cm,height=5.2cm,xlabel={\small year},ylabel={\small cumulative return (\%)},xmin=0,xmax=20,ymin=-20,ymax=420,
  tick label style={font=\footnotesize},grid=major,grid style={gray!20},table/col sep=comma,
  legend style={font=\scriptsize,at={(0.5,-0.3)},anchor=north,/tikz/every even column/.append style={column sep=8pt}},legend columns=2]
\addplot[omThm,thick] table[x=year,y=base]{figdata/desk/03-the-multi-manager-platform/overlay.csv};
\addplot[omDef,thick,dashed] table[x=year,y=hedged]{figdata/desk/03-the-multi-manager-platform/overlay.csv};
\legend{pods only,with the factor overlay}
\end{axis}
\end{tikzpicture}

\omcaption{The chapter's platform over twenty years (seed 1), gross of fees, with and without a centre-book overlay that hedges the pods' summed
exposure to their one shared factor. The overlay keeps the mean and removes the shared swings. Data: \texttt{fm\_platform.sim}, \texttt{firm.podshop.overlay}.}\label{fig:fm:the-multi-manager-platform:overlay}
\end{omfigure}

\begin{remark}[Who owns the centre book's P\&L]\label{rem:fm:the-multi-manager-platform:owner}
A \omterm{def:fm:the-multi-manager-platform:centre}{centre book}'s profit belongs to the fund, not to any team, so it earns no team payout; its costs pass through like any other. When the centre
adds to positions many pods hold, it increases the concentration the pods already create, and its risk belongs in the same limits as theirs
(chapter 7).
\end{remark}

\section{Tutorial: the investor's bill}\label{tut:fm:the-multi-manager-platform:tut}

\begin{tutorial}
\textbf{Goal.} Price a platform's fee structure and its \omterm{def:fm:the-multi-manager-platform:netting}{netting risk}. \textbf{End state:}
\cref{fig:fm:the-multi-manager-platform:waterfall,fig:fm:the-multi-manager-platform:netting}.
\begin{tutsteps}
\item \textbf{Pods.} \texttt{firm.podshop.pods(50, 20, 1.0, 0.20, 0.1, rng)} simulates daily pod returns; \texttt{annual} sums them by year.
\item \textbf{Payouts.} \texttt{payouts(ann, 0.20, carry=True, fire=0.10)} pays each team on its own profit
(\cref{lst:fm:the-multi-manager-platform:payouts}).
\item \textbf{The bill.} \texttt{fm\_platform.waterfall()} and \texttt{classic()} compare the pass-through structure with a 2-and-20 fund on the
same pods.
\item \textbf{Netting.} \texttt{fm\_platform.netting\_vs\_sr()} compares the simulation with \cref{prop:fm:the-multi-manager-platform:netting}.
\item \textbf{Ladder and overlay.} \texttt{ladder\_effect()} and \texttt{overlay\_effect()}.
\end{tutsteps}
\end{tutorial}

\textbf{What to change next.} Keep every team forever (exercise 7); correlate the pods at 0.4 and see the overlay's gain shrink once only part of
the shared risk is a tradable factor.

\section{Build: the platform model}\label{bld:fm:the-multi-manager-platform:podshop}

\begin{build}
\bfield{Purpose} The economics of a platform in one module: payouts, pass-through, the manager's allocation, \omterm{def:fm:the-multi-manager-platform:netting}{netting risk}, a drawdown ladder and a
centre-book overlay.
\bfield{Interface} \texttt{firm.podshop}: \texttt{pods(n, years, sr, vol, rho, rng, days, beta)}, \texttt{annual}; \texttt{payouts(ann, rate,
carry, fire)}, \texttt{netting\_cost}; \texttt{expected\_positive(mu, sigma)}; \texttt{Terms}, \texttt{waterfall(ann, terms, fire)},
\texttt{classic(ann, mgmt, perf)}; \texttt{ladder(daily, cut, stop)}; \texttt{overlay(daily, beta, factor)}.
\bfield{Rules} A team's payout never depends on another team's result; carried losses reduce only its own later payouts; a replaced team's carry is
forgotten; the waterfall adds up exactly.
\bfield{Acceptance tests} \texttt{code/firm/podshop/tests/}: carry-forward and no clawback on a hand example; the closed form against Monte
Carlo; pods' mean, volatility and correlation; the waterfall identity; the ladder cuts then stops; the overlay lowers fund volatility; replacement
restores payouts.
\bfield{Stretch} Allocate capital by trailing risk-adjusted return (firm.multistrat's \texttt{allocate}); a guaranteed first-year payout for new
teams; the \omterm{def:fm:the-multi-manager-platform:centre}{centre book} adding to the pods' consensus positions and its effect on concentration.
\end{build}

\begin{omsources}
\item Millennium Management LLC, Form ADV Part 2A brochure, 25 August 2026 (SEC Investment Adviser Public Disclosure).
\item One Quant Book 8, chapter 28 (pods, drawdown limits), and One Quant Book 7, chapter 28 (crowding).
\item S.~J.~Grossman and Z.~Zhou, ``Optimal investment strategies for controlling drawdowns'', \emph{Mathematical Finance} 3(3), 1993.
\end{omsources}

\section{Exercises}

\begin{exercise}[$\star$]\label{exo:fm:the-multi-manager-platform:1}
A team runs \$500 million of the fund's capital and makes 12\% in a year. At a \omterm{def:fm:the-multi-manager-platform:payout}{payout rate} of 20\%, what is it paid?
\end{exercise}

\begin{exercise}[$\star$]\label{exo:fm:the-multi-manager-platform:2}
The same team lost 5\% the year before and was kept. What is it paid on this year's 12\%, with losses carried forward?
\end{exercise}

\begin{exercise}[$\star$]\label{exo:fm:the-multi-manager-platform:3}
A pod's annual profit is normal with a Sharpe ratio of 1. Compute $\E[\max(X,0)]$ in units of its volatility.
\end{exercise}

\begin{exercise}[$\star\star$]\label{exo:fm:the-multi-manager-platform:4}
In \cref{ex:fm:the-multi-manager-platform:two}, each pod costs \$3 million a year to run. What does the investor pay in total, and what does the
fund earn for it?
\end{exercise}

\begin{exercise}[$\star\star$]\label{exo:fm:the-multi-manager-platform:5}
Use \cref{prop:fm:the-multi-manager-platform:netting} to compute the netting cost as a share of gross profit at a \omterm{def:fm:the-multi-manager-platform:payout}{payout rate} of 20\% for pods of
Sharpe ratio 0.5.
\end{exercise}

\begin{exercise}[$\star\star$]\label{exo:fm:the-multi-manager-platform:6}
The chapter's ladder cuts 88\% of team-years. Explain why from the pods' volatility, and propose thresholds that would cut about a quarter.
\end{exercise}

\begin{exercise}[$\star\star\star$]\label{exo:fm:the-multi-manager-platform:7}
\emph{Coding.} Rerun \texttt{fm\_platform.waterfall(fire=None)}: no team is ever replaced. What do team payouts and the investor's share become,
and what is the netting cost as a share of gross profit?
\end{exercise}

\begin{exercise}[$\star\star\star$]\label{exo:fm:the-multi-manager-platform:8}
\emph{Find the flaw.} ``Our teams' average Sharpe ratio is 1.5 over their track records, so \omterm{def:fm:the-multi-manager-platform:netting}{netting risk} costs us almost nothing.''
\end{exercise}

\section{Problem: Paid for the Winners}

\begin{problem}[{Weekend problem --- paid for the winners}]\label{pb:fm:the-multi-manager-platform:1}
An investor is offered a platform fund with fifty pods and a \omterm{def:fm:the-multi-manager-platform:pass}{pass-through fee} structure. You have its brochure and a simulator.

\medskip\noindent\textbf{Part I --- The structure.}
\begin{enumerate}
\item Define a \omterm{def:fm:the-multi-manager-platform:pass}{pass-through fee} and a \omterm{def:fm:the-multi-manager-platform:payout}{payout rate}.
\item From the brochure, how is a portfolio manager's compensation determined, and what happens to a team's losses?
\item What does the centre do with the fund's capital besides allocating it to teams?
\item Why can a platform's investors pay bonuses in a year the fund loses?
\end{enumerate}

\medskip\noindent\textbf{Part II --- The bill.}
\begin{enumerate}[resume]
\item Give the waterfall per \$100 of capital: gross, team payouts, costs, manager, investor.
\item In how many of the 400 simulated years did the fund lose?
\item What would a 2-and-20 fund leave its investors on the same pods, and why can it not exist?
\item What is the investor's share of gross profit?
\end{enumerate}

\medskip\noindent\textbf{Part III --- Netting.}
\begin{enumerate}[resume]
\item Define \omterm{def:fm:the-multi-manager-platform:netting}{netting risk} and state \cref{prop:fm:the-multi-manager-platform:netting}.
\item Compute the closed-form cost for Sharpe ratios 1, 0.5 and 0.25.
\item Give the simulated cost with teams replaced after a 10\% loss, at Sharpe ratios 1 and 0.5.
\item Give the simulated cost when no team is ever replaced, and explain the difference.
\item Which two properties of a platform set its \omterm{def:fm:the-multi-manager-platform:netting}{netting risk}?
\end{enumerate}

\medskip\noindent\textbf{Part IV --- The centre.}
\begin{enumerate}[resume]
\item Give the fund's volatility and Sharpe ratio with and without the \omterm{def:fm:the-multi-manager-platform:centre}{factor overlay}, and explain why the gain is an upper bound.
\item Give the shares of team-years cut and stopped by the ladder and its effect on return, volatility and Sharpe ratio.
\item Why does the ladder only cost in this model?
\item Define a \omterm{def:fm:the-multi-manager-platform:centre}{centre book} and say who owns its P\&L.
\item How should the \omterm{def:fm:the-multi-manager-platform:centre}{centre book}'s risk be limited?
\item State the \emph{named result}: the netting cost per pod in closed form, its value at Sharpe ratios 1 and 0.5, and the pass-through share of
gross profit (payouts plus costs) for fifty pods.
\item In two sentences, what should the investor ask the platform before investing?
\end{enumerate}
\end{problem}

\section{Interview questions}

\begin{interviewq}[$\star$ \iqroles{trader, researcher}]\label{iq:fm:the-multi-manager-platform:1}
What is a \omterm{def:fm:the-multi-manager-platform:pass}{pass-through fee}, and why do multi-manager platforms use it instead of a fixed management fee?
\end{interviewq}

\begin{interviewq}[$\star$ \iqroles{trader}]\label{iq:fm:the-multi-manager-platform:2}
You run a pod that lost 5\% last year and made 12\% this year. How is your payout computed, and what if you had been replaced?
\end{interviewq}

\begin{interviewq}[$\star\star$ \iqroles{researcher, risk}]\label{iq:fm:the-multi-manager-platform:3}
Two pods with the same Sharpe ratio: why does paying each on its own profit cost more than paying on their sum? Give the formula.
\end{interviewq}

\begin{interviewq}[$\star\star$ \iqroles{risk}]\label{iq:fm:the-multi-manager-platform:4}
Fifty pods with volatility 20\% each and pairwise correlation 0.1: what is the fund's volatility, and what if the shared part were hedged?
\end{interviewq}

\begin{interviewq}[$\star\star$ \iqroles{trader}]\label{iq:fm:the-multi-manager-platform:5}
Why is a 10\% drawdown cut a harsh rule for a pod with 20\% volatility? What would you propose?
\end{interviewq}

\begin{interviewq}[$\star\star\star$ \iqroles{researcher}]\label{iq:fm:the-multi-manager-platform:6}
A platform says its \omterm{def:fm:the-multi-manager-platform:centre}{centre book} only hedges. What in its P\&L and positions would you examine to test that claim?
\end{interviewq}
